\documentclass[10pt,aspectratio=169]{beamer}

\usepackage{spbu}

\usepackage[T2A]{fontenc}
\usepackage[utf8]{inputenc}
\usepackage[english,russian]{babel}
\usepackage{ulem}

\title{\LaTeX}
\subtitle{Набор математических формул}


\begin{document}

    \maketitle
    
    \begin{frame}[fragile]{Добавление формул}
    	Формулы можно разделить на два вида.
    	\begin{itemize}
    		\item Строчные формулы добавляются с использованием символа доллар (\verb|$|), например\\
    		\verb|$\Delta =\nabla^{2}$|: $\Delta =\nabla^{2}$.
    		\item Вынесенные формулы добавляются, например, с помощью двух символов доллара (\verb|$$формула$$|), либо с использованием конструкции \verb|\[формула\]|, либо с использованием окружения
    		\verb|equation| или \verb|displaymath|, существуют и другие способы.
    	\end{itemize}
    	
    	\begin{itemize}
    		\item Внутритекстовые и вынесенные формулы --- часть предложения.
    	\end{itemize}
    \end{frame}
    
    \begin{frame}[fragile]{Добавление формул}
    	\framesubtitle{Пример вынесенной формулы}
    	
    	\begin{block}{Добавление вынесенной формулы}
\begin{lstlisting}
\begin{equation}
	\frac{\partial u}{\partial t}
	-a^2\left(\frac{\partial^2 u}{\partial x^2_1}
	+\frac{\partial^2 u}{\partial x^2_2}
	+\ldots+\frac{\partial^2 u}{\partial x^2_n}\right)=f\left(x,t\right).
	\label{eq:eq1}
\end{equation}
\end{lstlisting}
    	\end{block}
    	
    	Результат:
    	\begin{equation}
    		\frac{\partial u}{\partial t}
    		-a^2\left(\frac{\partial^2 u}{\partial x^2_1}
    		+\frac{\partial^2 u}{\partial x^2_2}
    		+\ldots+\frac{\partial^2 u}{\partial x^2_n}\right)=f\left(x,t\right).
    		\label{eq:eq1}
    	\end{equation}
    	Ссылка на формулу (\ref{eq:eq1}).
    \end{frame}
    
    \begin{frame}[fragile]{Добавление формул}
    	\framesubtitle{Нумерация, пакет autonum}
    	По умолчанию формулы заданные через окружение \verb|equation| --- нумеруются, через окружение \verb|displaymath| --- нет.
    	
    	Для организации нумерации, только тех формул, на которые есть ссылка, можно использовать пакет  \verb|autonum|. Его достаточно только подключить.
    	
    	\begin{columns}
    		\begin{column}{0.5\textwidth}
\begin{lstlisting}
\begin{equation}
	a+b=2
	\label{eq:example1}
\end{equation}
Ссылка на формулу (2).
\end{lstlisting}
    		\end{column}
    		\begin{column}{0.5\textwidth}
\begin{lstlisting}
\begin{equation}
	a-b=0
	\label{eq:example2}
\end{equation}
Ссылка на формулу (\ref{eq:example2}).
\end{lstlisting}
    		\end{column}
    		% Может быть и больше!
    	\end{columns}
    \end{frame}
    
    \begin{frame}[fragile]{Пакеты}
       	\framesubtitle{amsmath, amsthm, amssymb, amsfonts}
       	
       	\begin{itemize}
    		\item Пакеты \verb|amsmath|, \verb|amsthm|, \verb|amssymb|, \verb|amsfonts|, разработаны AMS (American Mathematical Society).
    		\item Существенно упрощают набор формул.
    		\item Далее предполагаем, что эти пакеты подключены.
    	\end{itemize}
    \end{frame}
    
    \begin{frame}[fragile]{Арифметические операторы}
    	\begin{columns}
    		\begin{column}{0.5\textwidth}
		    	\begin{itemize}
    				\item \verb|$a+b$|: $a+b$
    				\item \verb|$a-b$|: $a-b$
    				\item \verb|$a \times b$|: $a \times b$
    				\item \verb|$a \cdot b$|: $a \cdot b$
    			\end{itemize}
    		\end{column}
    		\begin{column}{0.5\textwidth}
		    	\begin{itemize}
    				\item \verb|$a / b$|: $a / b$
		    		\item \verb|$a \div b$|: $a \div b$
    				\item \verb|$\frac{a+1}{b-5}$|: $\frac{a+1}{b-5}$
    				\item \verb|$\dfrac{a+1}{b-5}$|: $\dfrac{a+1}{b-5}$
    				\item \verb|$\tfrac{a+1}{b-5}$|: $\tfrac{a+1}{b-5}$
    			\end{itemize}
    		\end{column}
    	\end{columns}
    \end{frame}
    
    \begin{frame}[fragile]{Индексы}
		\begin{itemize}
			\item Надстрочные: $x^2$, $x^x$
		\end{itemize}
	\end{frame}
	
 	\begin{frame}[fragile]{Источники}
		\begin{itemize}
			\item Репозиторий на GitHub:
				\url{https://github.com/gleb-kun/mpt-course/tree/master/latex}
    		\item Кнут Д. Все про \TeX
			\item Львовский С. М. Набор и верстка в системе \LaTeX
			\item TikZ Manual for Version 2.0
			\begin{itemize}
				\item The TikZ and PGF Packages Manual for Version 3.1.10 \url{https://tikz.dev}
			\end{itemize}
			\item Воронцов К. В. \LaTeX $2{\varepsilon}$ в примерах
			\item Коробейников А. И. Презентации \url{https://statmod.ru/wiki/study:fall2018:tex}
			\item Лебедев А. Ководство
			\begin{itemize}
			\item \url{https://www.artlebedev.ru/kovodstvo/sections/97}
			\end{itemize}
			\item Мильчин А., Людмила Ч. Справочник издателя и автора
			\begin{itemize}
				\item \url{https://www.artlebedev.ru/izdal/spravochnik-izdatelya-i-avtora}
			\end{itemize}
		\end{itemize}
	\end{frame}

    \backmatter

\end{document}
